Suppose $G=\<\D\backslash\{e\}\mid
r_1,r_2,\dots\>$ is a presentation of $G$, i.e. the kernel of
the map induced by
$\sim$ from $F(\D\backslash\{e\})$ to $G$ is the normal closure of
the
$r_i$'s. Then each $r_i$ may be represented by a $k$-box, written
$\blankbox{r_i}$, with $\mu({\op{alt}}(\p(\blankbox{r_i})))=r_i$,
for some $k$ with $2k\geq\ell(r)$. (We use $\ell(w)$ to denote the
length of a word $w$.) To do this one may have to use $e\in\D$ so
that the word $r_i$ conforms with the alternating condition. For
instance to represent $\g\d^2\g^{-1}\d$ one might use the basic
tensor 
\fbox{\parbox{1.8cm}{
$\dsize{
\begin{array}{c}
\g\ e\ \d\ e\\
e\ \d\ \g\  \d 
\end{array}
        }$ 
}}.      

%%%%%%%%%%%%%%%%%%%%%%%%%%%%%%%%%%%%%%%%%%%%%%%
%$\boxed{\matrix \g & e & \d & e\\ e & \d &\g &
%\d\endmatrix}$.
%%%%%%%%%%%%%%%%%%%%%%%%%%%%%%%%%%%%%%%%%%%%%%%%%
Let $\mu:F(\D)\to F(\D\backslash\{e\})$ be the homomorphism defined
by $\mu(e)=1\in F(\D\backslash\{e\})$, \ $\mu|_{\D\backslash\{e\}}=$
id.

\bigskip{\bf Definition.} Let $R=\bigcup^{\infty}_{k=0} R_k$
be the planar subalgebra of $P^{\G,\D}$ generated by
$\{ \boxf{\d}  :\d\in\D\}\cup \{ \blankbox{r}\}
\cup \{\blankbox{\hskip2pt r^{-1}_i}\}$.
Let $H=\{\mu({\op{alt}}(\p T))\mid T$ is a basic tensor in $R\}$.

{\bf Theorem 2.7.6} The set $H$ is a subgroup of
$F(\D\backslash\{e\})$ equal to the normal closure $N$ of $\{r_i\}$
in $F(\D\backslash\{e\})$. Moreover, $P^{G,\D}=R$.

\medskip
{\bf Proof}. That $H$ is multiplicatively closed follows from
Figure 2.7.7.
%%%%%%%%%%%%%%%%%%%%%%%%%%%%%%%%%%%%%%%%%%%%%
\[
        \begin{picture}(0,0)%
\epsfig{file=ps/pic33.ps}%
\end{picture}%
\setlength{\unitlength}{0.00041700in}%
%
\begingroup\makeatletter\ifx\SetFigFont\undefined
% extract first six characters in \fmtname
\def\x#1#2#3#4#5#6#7\relax{\def\x{#1#2#3#4#5#6}}%
\expandafter\x\fmtname xxxxxx\relax \def\y{splain}%
\ifx\x\y   % LaTeX or SliTeX?
\gdef\SetFigFont#1#2#3{%
  \ifnum #1<17\tiny\else \ifnum #1<20\small\else
  \ifnum #1<24\normalsize\else \ifnum #1<29\large\else
  \ifnum #1<34\Large\else \ifnum #1<41\LARGE\else
     \huge\fi\fi\fi\fi\fi\fi
  \csname #3\endcsname}%
\else
\gdef\SetFigFont#1#2#3{\begingroup
  \count@#1\relax \ifnum 25<\count@\count@25\fi
  \def\x{\endgroup\@setsize\SetFigFont{#2pt}}%
  \expandafter\x
    \csname \romannumeral\the\count@ pt\expandafter\endcsname
    \csname @\romannumeral\the\count@ pt\endcsname
  \csname #3\endcsname}%
\fi
\fi\endgroup
\begin{picture}(5400,3044)(1201,-4284)
\put(3302,-1937){\makebox(0,0)[lb]{\smash{\SetFigFont{12}{14.4}{rm}$S$}}}
\put(3527,-3737){\makebox(0,0)[lb]{\smash{\SetFigFont{12}{14.4}{rm}$T$}}}
\put(1201,-2836){\makebox(0,0)[lb]{\smash{\SetFigFont{12}{14.4}{rm}$Q=$}}}
\put(6601,-2836){\makebox(0,0)[lb]{\smash{\SetFigFont{12}{14.4}{rm}alt($\partial Q)=$alt$(\partial S)$alt$(\partial T)$}}}
\end{picture}

\]
\begin{center}
Figure 2.7.7
\end{center}
%%%%%%%%%%%%%%%%%%%%%%%%%%%%%%%%%%%%%%%%%%%%%%
\noindent To see that $H=H^{-1}$, note that the transpose of a
basic tensor gives the inverse boundary word, and a planar algebra
generated by a $*$-closed set of boxes is $*$-closed (note that the
box $\blankbox{\d}$ is self-adjoint). Figure 2.7.8 exhibits
conjugation of $\alpha={\op{alt}}(\p T)$ by $\g\d\g^{-1}$ for
$\g,\d\in\D$, which shows how to prove that $H$ is normal
%%%%%%%%%%%%%%%%%%%%%%%%%%%%%%%%%%%%%%%%%%%
\[
        \begin{picture}(0,0)%
\epsfig{file=ps/pic34.ps}%
\end{picture}%
\setlength{\unitlength}{0.00041700in}%
%
\begingroup\makeatletter\ifx\SetFigFont\undefined
% extract first six characters in \fmtname
\def\x#1#2#3#4#5#6#7\relax{\def\x{#1#2#3#4#5#6}}%
\expandafter\x\fmtname xxxxxx\relax \def\y{splain}%
\ifx\x\y   % LaTeX or SliTeX?
\gdef\SetFigFont#1#2#3{%
  \ifnum #1<17\tiny\else \ifnum #1<20\small\else
  \ifnum #1<24\normalsize\else \ifnum #1<29\large\else
  \ifnum #1<34\Large\else \ifnum #1<41\LARGE\else
     \huge\fi\fi\fi\fi\fi\fi
  \csname #3\endcsname}%
\else
\gdef\SetFigFont#1#2#3{\begingroup
  \count@#1\relax \ifnum 25<\count@\count@25\fi
  \def\x{\endgroup\@setsize\SetFigFont{#2pt}}%
  \expandafter\x
    \csname \romannumeral\the\count@ pt\expandafter\endcsname
    \csname @\romannumeral\the\count@ pt\endcsname
  \csname #3\endcsname}%
\fi
\fi\endgroup
\begin{picture}(5415,2444)(1201,-3983)
\put(1201,-2836){\makebox(0,0)[lb]{\smash{\SetFigFont{12}{14.4}{rm}$Q=$}}}
\put(2551,-2836){\makebox(0,0)[lb]{\smash{\SetFigFont{12}{14.4}{rm}$\gamma$}}}
\put(3181,-2866){\makebox(0,0)[lb]{\smash{\SetFigFont{12}{14.4}{rm}$e$}}}
\put(3736,-2851){\makebox(0,0)[lb]{\smash{\SetFigFont{12}{14.4}{rm}$\delta$}}}
\put(4321,-2851){\makebox(0,0)[lb]{\smash{\SetFigFont{12}{14.4}{rm}$\gamma$}}}
\put(5281,-2836){\makebox(0,0)[lb]{\smash{\SetFigFont{12}{14.4}{rm}$T$}}}
\put(6616,-3361){\makebox(0,0)[lb]{\smash{\SetFigFont{12}{14.4}{rm}$=\gamma\delta\gamma^{-1}$alt$(\partial Q)\gamma\delta^{-1}\gamma^{-1}$}}}
\put(6601,-2836){\makebox(0,0)[lb]{\smash{\SetFigFont{12}{14.4}{rm}$\mu$(alt($\partial T)$}}}
\end{picture}

\]
\begin{center}
Figure 2.7.8
\end{center}
%%%%%%%%%%%%%%%%%%%%%%%%%%%%%%%%%%%%%%%%%%%%%%%%%%%%
\noindent Thus $H$ contains the normal closure $N$.

Now the tangle picture of an arbitrary basic tensor in $R$ can be
isotoped so that it is as in Figure 2.7.9.
\[
        \begin{picture}(0,0)%
\epsfig{file=ps/pic35.ps}%
\end{picture}%
\setlength{\unitlength}{0.00041700in}%
%
\begingroup\makeatletter\ifx\SetFigFont\undefined
% extract first six characters in \fmtname
\def\x#1#2#3#4#5#6#7\relax{\def\x{#1#2#3#4#5#6}}%
\expandafter\x\fmtname xxxxxx\relax \def\y{splain}%
\ifx\x\y   % LaTeX or SliTeX?
\gdef\SetFigFont#1#2#3{%
  \ifnum #1<17\tiny\else \ifnum #1<20\small\else
  \ifnum #1<24\normalsize\else \ifnum #1<29\large\else
  \ifnum #1<34\Large\else \ifnum #1<41\LARGE\else
     \huge\fi\fi\fi\fi\fi\fi
  \csname #3\endcsname}%
\else
\gdef\SetFigFont#1#2#3{\begingroup
  \count@#1\relax \ifnum 25<\count@\count@25\fi
  \def\x{\endgroup\@setsize\SetFigFont{#2pt}}%
  \expandafter\x
    \csname \romannumeral\the\count@ pt\expandafter\endcsname
    \csname @\romannumeral\the\count@ pt\endcsname
  \csname #3\endcsname}%
\fi
\fi\endgroup
\begin{picture}(9029,2722)(894,-3383)
\end{picture}

\]
\begin{center}
        Figure 2.7.9
\end{center}

\noindent This is a tangle $T$ all of whose curves are vertical
straight lines surrounded by an element of the category ${\cal A} (\emptyset)$. But it
is easy to see that applying an ${\cal A} (\emptyset)$ element changes alt$(\p T)$ at
most by a conjugation.

The most difficult part of Theorem 2.7.6 is to show that
$P^{\G,\D}=R$. We must show that if $w$ is a word of even length on
$\D$ with $\phi({\op{alt}}(w))=1$ then there is a basic tensor
$T\in R$ with $\p T=w$.

As a first step, observe that if $\mu$ is a homomorphism from
$F(\D)$ to $F(\D-\{e\})$ sending $e$ to the identity and with
$\mu(\d)=\d$ for $\d\neq e$, then if $w_1,w_2\in\frak S(\D)$ are of
even length and $\mu({\op{alt}}(w_1))=\mu({\op{alt}}(w_2))$ then
${\op{alt}}(w_1)={\op{alt}}(w_2)$. This is because
$w_1w^*_2$ ($w^*$ is $w$ written backwards) satisfies
${\op{alt}}(w_1w^*_2)={\op{alt}}(w_1){\op{alt}}(w_2)^{-1}$, thus
${\op{alt}}(w_1w_2^*)\in{\op{ker}} \ \mu$ which is the normal
closure of $e$. The length of $w_1w_2^*$ can be reduced
(if necessary) by eliminating consecutive letters two at a time to
obtain another word $w$, of even length, with
$\ell(w)=\ell({\op{alt}} \ w)$ ($\ell=$ length). By the uniqueness
of reduced words in a free group, $w$ must be a product of words of
the form $x \ e \ y$, which map to conjugates of
$e^{\pm 1}$ in $F(\D)$. But the last letter of $x$ and the first
letter of $y$ must then be the same, and alt will send both these
letters to the same free group element. Thus in the process of
reducing $w_1w^*_2$, all occurrences of $e$ must disappear and
${\op{alt}}(w_1)={\op{alt}}(w_2)$.

A consequence of this observation is that, if $T$ is a basic tensor
with $\phi({\op{alt}}(\p T))=1$ in $\G$ so that
$\mu({\op{alt}}(\p T))$ is in the normal closure of $\{r_i\}$ in
$F(\D\backslash\{e\})$, then alt$(\p T)$ is the normal closure of
$\{{\op{alt}}(\p( \blankbox{r_i}))\}$ in $F(\D)$. Thus it
suffices to show that, if $T_1$ and $T_2$ are basic tensors with
${\op{alt}}(\p T_1)={\op{alt}}(\p T_2)\in F(\D)$, then $T_1=cT_2$
for some $c$ in ${\cal A} (\emptyset)$ (since for any $x\in\frak
S(\D)$ with $\phi |\mu({\op{alt}}(x))=1$ we have shown there is a
$T$ in the planar algebra $R$ with alt$(\p T)={\op{alt}}(x)$).
But this is rather easy --- we may suppose without loss of
generality that no cancellation happens going from $\p T_1$ to
alt$(\p T_1)$ and then use induction on $\ell(\p T_2)$. If
$\ell(\p T_2)=\ell(\p T_1)$ then $\p T_2=\p T_1$. Otherwise there
must be a sequence $\dots \d\d\dots$ in $\p T_2$ for some
$\d\in\D$. Connecting $\d$ to $\d$ in the tangle reduces the length
of $\p T_2$ by 2, and the remaining region is a disc.

\qed
A.Casson has pointed out the connection between
$P^{\G,\D}$ and van Kampen diagrams.

\bigskip
If $e\in\D$, the dimension of $P^{\G,\D}_n$ is the number of ways
of writing $1\in\G$ as a product of elements $\tilde\d$,
$\d\in\D$, with alternating signs. In particular,
dim$(P^{\G,\D}_1)=|\D|^2$ iff $\G$ is trivial. Since the problem of
the triviality of a group with given presentation is undecidable,
we conclude the following.

\bs\ni{\bf Corollary 2.7.10} The calculation of the dimension of a
planar subalgebra of $P^{\otimes}$ is undecidable.

\bs
Since there are
groups which are finitely generated but not finitely presented we
have

\bs\ni{\bf Corollary 2.7.11} There are finite dimensional
planar $*$-algebras which are not finitely generated as planar
algebras.

\bs
{\sc Proof}. If finitely many linear combinations of basic tensors
generated a planar algebra, then certainly the basic tensors
involved would also. But by 2.7.6, the group would then be finitely
presented. \qed

\bs{\bf Example 2.8 \ Spin models.}
We give a general planar algebra that is not a planar algebra, although it
is the planar algebras associated with it that will be of most
interest. In some sense it is a ``square root" of the planar
algebra $P^{\otimes}$ of $\S$2.6.

Let $V$ be a vector space of dimension $Q$ with a basis indexed by
``spin states" $\{1,2,\dots Q\}$. For each odd $n$ let $P^\s_n$ be
the subalgebra of End$(V^{\otimes\frac{n+1}{2}})$ given by
End$(V^{\otimes\frac{n-1}{2}})\otimes\D$ where $\D$ is the
subalgebra of End$(V)$ consisting of linear maps diagonal with
respect to the basis. For $n=0$, $P_0$ is the field $K$ and for $n$
even, $P^\s_n={\op{End}}(V^{\otimes\frac{n}{2}})$. Elements of
$P^\s_n$ will be identified with functions from $\{1,2,\dots ,Q\}^n$
to $K$, the value of the function on $(i_1,i_2,\dots ,i_n)$ being
the coefficient of the basic tensor
\fbox{\parbox{4.3cm}{
$\dsize{
\begin{array}{clcr}
i_1 & i_2  \ldots  \ldots i_{m-1} & i_m \\
i_n & i_{n-1}  \ldots .i_{m+2} & i_{m+1} 
\end{array}
        }$ 
}}      

%%%%%%%%%%%%%%%%%%%%%%%%%%%%%%%%%%%%%%%%%%%%%%%%%%%%%%%%%%%%%%%
%$\boxed{\matrix\format\l&\l&\r&\r&\r\\ i_1 \quad i_2\qquad \dots
%\quad i_{m-1}\quad i_m\\ i_n \quad i_{n-1}\quad \dots \quad
%i_{m+2}\quad i_{m+1}\endmatrix}$ 
%%%%%%%%%%%%%%%%%%%%%%%%%%%%%%%%%%%%%%%%%%%%%%%%%%%%%%%%%%%%%
for
$n=2m$, and the coefficient of
\fbox{\parbox{4.3cm}{
$\dsize{
\begin{array}{clcr}
i_1  & i_2  \ldots \ldots i_{m-1} & i_m \\
i_n &  i_{n-1}  \ldots .i_{m+2} & i_{m}      
\end{array}
        }$
}}

%%%%%%%%%%%%%%%%%%%%%%%%%%%%%%%%%%%%%%%%%%%%%%%%%%%%%%%%%%%%%%%%
%$\boxed{\matrix\format\l&\l&\r&\r&\r\\ i_1 \quad i_2\qquad \dots
%\quad i_{m-1}\quad i_m\\ i_n \quad i_{n-1}\quad \dots \quad
%i_{m+1}\quad i_m\endmatrix}$
%%%%%%%%%%%%%%%%%%%%%%%%%%%%%%%%%%%%%%%%%%%%%%%%%%%%%%%%%%%%%%%%%
for
$n=2m-1$. (See $\S$2.6 for notation.) We shall make $P^\s$ into a
general planar algebra in two slightly different ways. In both cases the
labelling set will be $P^\s$ itself.

\bs{\bf First planar structure on $P^\s$.}
Take a tangle $T$ in ${\cal P}_k(L)$. We will define
$\Phi_0(T)\in{\cal P}^\s_k$ as follows.

First, shade the connected components of ${\cal B}_k\backslash T$
(called regions) black and white so that the region containing a
neighborhood of (0,0) is white, and so that regions whose closures
intersect (i.e.~which share an edge) have different colours. In
other words, regions whose boundary induces the positive
orientation of $\Bbb R^2$ are coloured white and negatively oriented
ones are black. Observe that the top and bottom of ${\cal B}_k$
consists of segments of length one forming parts of the boundaries
of regions alternately coloured white and black. If $k$ is odd, the
right-hand boundary of ${\cal B}_k$ can be joined with the rightmost
top and bottom segments to form part of the boundary of a black
region.  This way the boundary of ${\cal B}_k$ always has $k$
segments attached to black regions whose closure meets the
boundary. Number these segments cyclically $1,2,\dots ,k$ starting
form the top left and going clockwise.  To define an element of
$P^\s_k$ from $T$ we must give a function 
$\Phi_0(T):\{1,2,\dots Q\}^k\to K$. It is
\[
\Phi_0(T)(i_1,i_2,\dots i_k)=\sum_\s \
\prod_{B\in\left\{
{\text{labelled\  boxes}}\atop  {\text{of }} T\right\}} \s (B)
\]
where $\s$ runs over all functions from the black regions of $T$ to
$\{1,2,\dots ,Q\}$ which take the value $i_j$ on the black region
whose closure contains the $j^{\text{th}}$ boundary segment, for
all $j=1,2,\dots ,k$. Given a labeled box $B$ of $T$, and such a
$\s$, the boundary segments of $B$ which meet closures of black
regions are numbered 1 to $k_B$ so $\s$ defines an element of
$\{1,2,\dots ,Q\}^{k_B}$, and thus the label of $B$ gives a scalar
$\s(B)$ in $K$. As usual empty sums are zero and empty products are
1. This completes the definition of $\Phi_0$ and it is easily
checked that $\Phi_0$ presents $P^\s$ as a planar algebra. The
induced representation of ${\cal P}(\phi)$ gives a representation of
TL with $\d_1\!=\!Q$, $\d_2\!=\!1$. It is precisely the
representation associated with the Potts model used by Temperley
and Lieb in [TL].

 To be sure of relevance to subfactors, we
now show how to adjust these parameters so that
$\d_1=\d_2=\sqrt{Q}$.

\bs{\bf Second planar structure on $P^\s$.}
If $T$ is a labeled tangle in ${\cal P}_k(L)$, we define
a tangle $\tilde T$ in ${\cal P}_k(\phi)$ by ``smoothing" all the
boxes of $\tilde T$, i.e.~replacing $\boxe{R}$ by
$\uparrow \ \downarrow \ \uparrow$, and shrinking all
non-through-strings to semicircles near the top or bottom of ${\cal B}_k$. 
Put $f(T)=Q^{\frac 12(n_+-n_-)+\frac 14(n^\p_+-n^\p_-)}$
where $n_+$ and $n_-$ are the numbers of positively and negatively
oriented circles in $\tilde T$ respectively, and $n_{\pm}^\p$ are
similarly the numbers of positively and negatively oriented
semicircles near the top and bottom.  Thus defined, $f(T)$ is
clearly an isotopy invariant, so we could redefine it by assuming
all the boxes are parallel to the $x$-axis. Assuming all maxima and
minima of the $y$-coordinate restricted to the strings of $T$ are
nondegenerate, $2(n_+-n_-)+(n^\p_+-n^\p_-)$ is just
$p_++q_+-p_--q_-$ where $p_+,p_-$ are the numbers of local maxima
of $y$ oriented to the left and right respectively and similarly
$q_+$ and $q_-$ count minima to the right and left respectively.
It follows that $T\to f(T)$ is multiplicative and indeed that if $A$
is in the annular category one may define $f(A)$ so that
$f(AT)=f(A)f(T)$.
%%%%%%%%%%%%%%%%%%%%% The following is the hand written part %%%%%%%%%%%%%%%%%%%


The normalisation constant
$n_+-n_-+\frac{1}{2}(n_+^{\partial}-n_-^{\partial})$
may seem mysteious. What is actually being calculated is
the isotopy invariant $\int d\theta$
where the intgral is taken over the strings of the tangle
and $d\theta$ is the change of angle or curvature 1-form, normalised so that
integrating over a positively oriented circle counts one.
The above factor is then this integral when all strings meet
all boxes at right angles. Thus, by following shaded regions at every internal
box, another formula for this normalisation factor for a $k-$tangle is

\[
        [\frac{k+1}{2}]-b-\sum_{i\geq 1}[\frac{n_i}{2}].
\]

If all boxes are 2- or 3-boxes we get
\[
        \frac{1}{2}(\#(\mbox{black regions})-\#(\mbox{boxes}))+
        \frac{1}{4}(n_+^{\partial}-n_-^{\partial}),
\]
where now $n_+^{\partial}$ and $n_-^{\partial}$ are calculated
by eliminating the boxes by following the black regions rather
than going straight through the box.



%%%%%%%%%%%%%%%%%%%%%%%%%%%%%%%%%%%%%%%%%%%%%%%%%%%%%%%%%%%%%%%%%%%%%%%%%%%
\bs
\ni{\bf Proposition 2.8.1} The map $\Phi^\s:{\cal P}(L)\to P^\s$,
$\Phi(T)=f(T)\Phi_0(T)$ (linearly extended) presents $P^\s$ as a
general planar algebra, $\Phi |_{{\cal P}(\phi)}$  presents $TL$ with
$\d_1=\d_2=\sqrt{Q}$.

\bs
{\sc Proof.} Annular invariance follows from the relation
$f(AT)=f(A)f(T)$ and the annular invariance of $\Phi_0$.
If $\deye$ is a part of a tangle $T$ then
$\Phi(T)=QQ^{-\frac 12} \Phi(\tilde T)$ where $\deye$
has been removed from $T$. If $\deyer$\ \ \  is part of $T$,
$\Phi(T)=Q^{\frac 12} \Phi(\tilde T)$. \qed

\bs
When we refer to $P^\s$, we will mean $P^\s$,
together with $\Phi$.

Although dim $P^\s_0=1$, so that $P^\s$ gives an invariant of
labeled planar networks with unbounded region positively oriented,
dim$(P^\s_{1,1})=Q$ so $P^\s$ is not planar. However, $P^\s$
does have the obvious $*$ structure and tr${}_R$ is defined and
positive definite, so that any self-adjoint planar
subalgebra of
$P^\s$ will be a $C^*$-planar algebra.

\bs\ni
{\bf Proposition 2.8.2} A planar subalgebra $P$ of
$P^\s$ is spherical.

\bs
{\sc Proof.} Given a planar network $N$ in $P$ with positively
oriented unbounded region, we need only show that
\[
        \begin{picture}(0,0)%
\epsfig{file=ps/pic36.ps}%
\end{picture}%
\setlength{\unitlength}{0.00041700in}%
%
\begingroup\makeatletter\ifx\SetFigFont\undefined
% extract first six characters in \fmtname
\def\x#1#2#3#4#5#6#7\relax{\def\x{#1#2#3#4#5#6}}%
\expandafter\x\fmtname xxxxxx\relax \def\y{splain}%
\ifx\x\y   % LaTeX or SliTeX?
\gdef\SetFigFont#1#2#3{%
  \ifnum #1<17\tiny\else \ifnum #1<20\small\else
  \ifnum #1<24\normalsize\else \ifnum #1<29\large\else
  \ifnum #1<34\Large\else \ifnum #1<41\LARGE\else
     \huge\fi\fi\fi\fi\fi\fi
  \csname #3\endcsname}%
\else
\gdef\SetFigFont#1#2#3{\begingroup
  \count@#1\relax \ifnum 25<\count@\count@25\fi
  \def\x{\endgroup\@setsize\SetFigFont{#2pt}}%
  \expandafter\x
    \csname \romannumeral\the\count@ pt\expandafter\endcsname
    \csname @\romannumeral\the\count@ pt\endcsname
  \csname #3\endcsname}%
\fi
\fi\endgroup
\begin{picture}(5578,1341)(1175,-1455)
\put(1917,-901){\makebox(0,0)[lb]{\smash{\SetFigFont{12}{14.4}{rm}$N$}}}
\put(3207,-901){\makebox(0,0)[lb]{\smash{\SetFigFont{12}{14.4}{rm}$=$}}}
\put(4032,-931){\makebox(0,0)[lb]{\smash{\SetFigFont{12}{14.4}{rm}$Z(N)$}}}
\end{picture}

\]
\ni But since $P$ is planar, the sum over all internal spins in
\begin{picture}(0,0)%
\epsfig{file=ps/pic37.ps}%
\end{picture}%
\setlength{\unitlength}{0.00041700in}%
%
\begingroup\makeatletter\ifx\SetFigFont\undefined
% extract first six characters in \fmtname
\def\x#1#2#3#4#5#6#7\relax{\def\x{#1#2#3#4#5#6}}%
\expandafter\x\fmtname xxxxxx\relax \def\y{splain}%
\ifx\x\y   % LaTeX or SliTeX?
\gdef\SetFigFont#1#2#3{%
  \ifnum #1<17\tiny\else \ifnum #1<20\small\else
  \ifnum #1<24\normalsize\else \ifnum #1<29\large\else
  \ifnum #1<34\Large\else \ifnum #1<41\LARGE\else
     \huge\fi\fi\fi\fi\fi\fi
  \csname #3\endcsname}%
\else
\gdef\SetFigFont#1#2#3{\begingroup
  \count@#1\relax \ifnum 25<\count@\count@25\fi
  \def\x{\endgroup\@setsize\SetFigFont{#2pt}}%
  \expandafter\x
    \csname \romannumeral\the\count@ pt\expandafter\endcsname
    \csname @\romannumeral\the\count@ pt\endcsname
  \csname #3\endcsname}%
\fi
\fi\endgroup
\begin{picture}(1543,1191)(1214,-1327)
\put(1842,-811){\makebox(0,0)[lb]{\smash{\SetFigFont{12}{14.4}{rm}$N$}}}
\end{picture}

is independent of the spin value in the unbounded
region, and each term in the sum for 
\begin{picture}(0,0)%
\epsfig{file=ps/pic37.ps}%
\end{picture}%
\setlength{\unitlength}{0.00041700in}%
%
\begingroup\makeatletter\ifx\SetFigFont\undefined
% extract first six characters in \fmtname
\def\x#1#2#3#4#5#6#7\relax{\def\x{#1#2#3#4#5#6}}%
\expandafter\x\fmtname xxxxxx\relax \def\y{splain}%
\ifx\x\y   % LaTeX or SliTeX?
\gdef\SetFigFont#1#2#3{%
  \ifnum #1<17\tiny\else \ifnum #1<20\small\else
  \ifnum #1<24\normalsize\else \ifnum #1<29\large\else
  \ifnum #1<34\Large\else \ifnum #1<41\LARGE\else
     \huge\fi\fi\fi\fi\fi\fi
  \csname #3\endcsname}%
\else
\gdef\SetFigFont#1#2#3{\begingroup
  \count@#1\relax \ifnum 25<\count@\count@25\fi
  \def\x{\endgroup\@setsize\SetFigFont{#2pt}}%
  \expandafter\x
    \csname \romannumeral\the\count@ pt\expandafter\endcsname
    \csname @\romannumeral\the\count@ pt\endcsname
  \csname #3\endcsname}%
\fi
\fi\endgroup
\begin{picture}(1543,1191)(1214,-1327)
\put(1842,-811){\makebox(0,0)[lb]{\smash{\SetFigFont{12}{14.4}{rm}$N$}}}
\end{picture}

is $Q$ times the
corresponding term for
\begin{picture}(0,0)%
\epsfig{file=ps/pic38.ps}%
\end{picture}%
\setlength{\unitlength}{0.00041700in}%
%
\begingroup\makeatletter\ifx\SetFigFont\undefined
% extract first six characters in \fmtname
\def\x#1#2#3#4#5#6#7\relax{\def\x{#1#2#3#4#5#6}}%
\expandafter\x\fmtname xxxxxx\relax \def\y{splain}%
\ifx\x\y   % LaTeX or SliTeX?
\gdef\SetFigFont#1#2#3{%
  \ifnum #1<17\tiny\else \ifnum #1<20\small\else
  \ifnum #1<24\normalsize\else \ifnum #1<29\large\else
  \ifnum #1<34\Large\else \ifnum #1<41\LARGE\else
     \huge\fi\fi\fi\fi\fi\fi
  \csname #3\endcsname}%
\else
\gdef\SetFigFont#1#2#3{\begingroup
  \count@#1\relax \ifnum 25<\count@\count@25\fi
  \def\x{\endgroup\@setsize\SetFigFont{#2pt}}%
  \expandafter\x
    \csname \romannumeral\the\count@ pt\expandafter\endcsname
    \csname @\romannumeral\the\count@ pt\endcsname
  \csname #3\endcsname}%
\fi
\fi\endgroup
\begin{picture}(713,945)(1587,-1029)
\put(1842,-811){\makebox(0,0)[lb]{\smash{\SetFigFont{12}{14.4}{rm}$N$}}}
\end{picture}
. Taking the sum over all $Q$
spin states in the shaded region we are done. \qed

\bs There are ways to obtain planar subalgebras of
$P^\s$. An obvious place to look is association schemes 
where one is given a family of (0,1) $Q\times Q$ matrices $A_i$,
$i=1,\dots ,d$, whose linear span is closed under the operations
given by the tangles
\medskip
%%%%%%%%%%%%%%%%%%%%%%%%%%%%%%%%%%%%%
\begin{picture}(0,0)%
\epsfig{file=ps/pic39.ps}%
\end{picture}%
\setlength{\unitlength}{0.00041700in}%
%
\begingroup\makeatletter\ifx\SetFigFont\undefined
% extract first six characters in \fmtname
\def\x#1#2#3#4#5#6#7\relax{\def\x{#1#2#3#4#5#6}}%
\expandafter\x\fmtname xxxxxx\relax \def\y{splain}%
\ifx\x\y   % LaTeX or SliTeX?
\gdef\SetFigFont#1#2#3{%
  \ifnum #1<17\tiny\else \ifnum #1<20\small\else
  \ifnum #1<24\normalsize\else \ifnum #1<29\large\else
  \ifnum #1<34\Large\else \ifnum #1<41\LARGE\else
     \huge\fi\fi\fi\fi\fi\fi
  \csname #3\endcsname}%
\else
\gdef\SetFigFont#1#2#3{\begingroup
  \count@#1\relax \ifnum 25<\count@\count@25\fi
  \def\x{\endgroup\@setsize\SetFigFont{#2pt}}%
  \expandafter\x
    \csname \romannumeral\the\count@ pt\expandafter\endcsname
    \csname @\romannumeral\the\count@ pt\endcsname
  \csname #3\endcsname}%
\fi
\fi\endgroup
\begin{picture}(644,2744)(1179,-2933)
\end{picture}

(matrix multiplication) and
\begin{picture}(0,0)%
\epsfig{file=ps/pic40.ps}%
\end{picture}%
\setlength{\unitlength}{0.00041700in}%
%
\begingroup\makeatletter\ifx\SetFigFont\undefined
% extract first six characters in \fmtname
\def\x#1#2#3#4#5#6#7\relax{\def\x{#1#2#3#4#5#6}}%
\expandafter\x\fmtname xxxxxx\relax \def\y{splain}%
\ifx\x\y   % LaTeX or SliTeX?
\gdef\SetFigFont#1#2#3{%
  \ifnum #1<17\tiny\else \ifnum #1<20\small\else
  \ifnum #1<24\normalsize\else \ifnum #1<29\large\else
  \ifnum #1<34\Large\else \ifnum #1<41\LARGE\else
     \huge\fi\fi\fi\fi\fi\fi
  \csname #3\endcsname}%
\else
\gdef\SetFigFont#1#2#3{\begingroup
  \count@#1\relax \ifnum 25<\count@\count@25\fi
  \def\x{\endgroup\@setsize\SetFigFont{#2pt}}%
  \expandafter\x
    \csname \romannumeral\the\count@ pt\expandafter\endcsname
    \csname @\romannumeral\the\count@ pt\endcsname
  \csname #3\endcsname}%
\fi
\fi\endgroup
\begin{picture}(1845,1618)(1179,-2107)
\end{picture}

(Hadamard product).
%%%%%%%%%%%%%%%%%%%%%%%%%%%%%%%%%%%%
\medskip\ni
The requirement that this linear span (the
``Bose-Mesner algebra") be closed under all planar contractions is
presumably much more stringent. If the requirement is satisfied, and the row and
column sums of each $A_i$ do not depend on the row or column, we
will have a planar subalgebra of
$P^\s$. A particularly simple example of this comes from transitive
actions of a finite group $G$ on a set $S$. Then the orbits of $G$
on $S\times S$ define an association scheme whose Bose-Mesner
algebra is the fixed points for the action on $M_{|S|}(\Bbb C)$ by
conjugation. We get a planar subalgebra of $P^\s$ either
by taking the fixed points for the $G$-action on $P^\s$ or the
planar subalgebra generated by the assocation scheme. They are
different in general. A case where they are the same is for the
dihedral group on a set with five elements (see [J4]). They
are different for Jaeger's Higman-Sims model ([Ja],[dlH])
 --- although the dimensions of the two planar algebras agree for
a while, they have different asymptotic growth rates, one being that
of the commutant of Sp(4) on $(\Bbb C^4)^{\otimes k}$ and the other
being $100^k$. 

Here is an interesting example for a doubly transitive group.
It connects with Example 2.5 and gives a new kind of ``spin model"
for link invariants from links projected with only triple point
singularities.

The alternating group $A_4$ is doubly transitive on the set
$\{1,2,3,4\}$ but there are two orbits on the set of ordered
triples $(a,b,c)$ of distinct elements according to whether
$1\mapsto a$, $2\mapsto b$, $3\mapsto c$, $4\mapsto d$
(with $\{a,b,c,d\}=\{1,2,3,4\}$) is an even or odd permutation.
Let \ $\boxe{e}\in{\cal P}_3(L)$ be such that
$\Phi^\s(\boxe{e})$ is the characteristic
function  of the even orbit. Define a mapping from
${\cal P}(\ \crsst\ )$ (the universal planar algebra on a single 3-box)
to $P^\s$ by sending $\crsst$ to
$\boxe{e}-\frac 12 \vln\ \  \smth$.
It is possible to prove that this map passes to the quotient
$P^H$ (the planar algebra of 2.5) with parameters  $t=i=x$.
This is equivalent to showing that twice the value of the {\sc
homfly} polynomial of a link obtained by connecting 3-boxes
$\crsst$ (at 1,--1 in $\ell-m$ variables) in an oriented way is
the partition function in $P^\s$ (with $Q=4$) given by filling the
same three boxes with $\boxe{e}-\frac 12 \vln\ \ \smth$. 
We give a sample calculation below
which illustrates all the considerations. Note that, for $t=i=x$,
the value of a single circle in the {\sc homfly} skein is 2.
\[
        \begin{picture}(0,0)%
\epsfig{file=ps/pic41.ps}%
\end{picture}%
\setlength{\unitlength}{0.00041700in}%
%
\begingroup\makeatletter\ifx\SetFigFont\undefined
% extract first six characters in \fmtname
\def\x#1#2#3#4#5#6#7\relax{\def\x{#1#2#3#4#5#6}}%
\expandafter\x\fmtname xxxxxx\relax \def\y{splain}%
\ifx\x\y   % LaTeX or SliTeX?
\gdef\SetFigFont#1#2#3{%
  \ifnum #1<17\tiny\else \ifnum #1<20\small\else
  \ifnum #1<24\normalsize\else \ifnum #1<29\large\else
  \ifnum #1<34\Large\else \ifnum #1<41\LARGE\else
     \huge\fi\fi\fi\fi\fi\fi
  \csname #3\endcsname}%
\else
\gdef\SetFigFont#1#2#3{\begingroup
  \count@#1\relax \ifnum 25<\count@\count@25\fi
  \def\x{\endgroup\@setsize\SetFigFont{#2pt}}%
  \expandafter\x
    \csname \romannumeral\the\count@ pt\expandafter\endcsname
    \csname @\romannumeral\the\count@ pt\endcsname
  \csname #3\endcsname}%
\fi
\fi\endgroup
\begin{picture}(7277,6618)(88,-5970)
\put(676,-3811){\makebox(0,0)[lb]{\smash{\SetFigFont{12}{14.4}{rm}$b$}}}
\put(2881,-3676){\makebox(0,0)[lb]{\smash{\SetFigFont{12}{14.4}{rm}$c$}}}
\put(3166,-4411){\makebox(0,0)[lb]{\smash{\SetFigFont{12}{14.4}{rm}$c$}}}
\put(5086,-5266){\makebox(0,0)[lb]{\smash{\SetFigFont{12}{14.4}{rm}$a$}}}
\put(4846,-2701){\makebox(0,0)[lb]{\smash{\SetFigFont{12}{14.4}{rm}$c$}}}
\put(2461,-1696){\makebox(0,0)[lb]{\smash{\SetFigFont{12}{14.4}{rm}$d$}}}
\end{picture}

\]
\begin{center}
        Figure 2.8.3
\end{center}

Smoothing all the 3-boxes leads to a single negatively oriented
circle so we must divide the final partition function by 2.
Replacing the 3-boxes by $\boxe{e}-\frac12 \vln\ \  \smth$ we look for spin states, i.e. functions
from the shaded regions to $\{1,2,3,4\}$ for which each 3-box
yields a non-zero contribution to the partition function. Around
each 3-box this means that either the three spin values are in the
even orbit under
$A_4$, or they are all the same. The first case contributes +1  to
the product over boxes, the second case contributes --1 ({\bf not}
$-\frac 12$ because of the maxima and minima in the box). If the box labeled
($\dagger$) is surrounded by the same spin value, all the spin
states must be the same for a nonzero contribution to $Z$. This
gives a factor $4\times (-1)^5$. On the other hand, if the spins at
($\dagger$) are as in Figure 2.8.3 with $(a,b,c)$ in the even
orbit, the other spin
choices are forced (where $\{a,b,c,d\}=\{1,2,3,4\}$), for a
contribution of --1. The orbit is of size 12 so the partition
function is $\frac 12(-12-4)=-8$. For this link the value of the
{\sc homfly} polynomial $P_L(1,-1)$ is --4. The factor of 2 is
accounted for by the fact that our partition function is 2 on the
unknot. Thus our answer is correct. Note how few spin patterns
actually contributed to $Z$!

If we wanted to use non-alternating 3-boxes we could simply use the
{\sc homfly} skein relation to modify the 3-box. For instance
\[
        \begin{picture}(0,0)%
\epsfig{file=ps/pic42.ps}%
\end{picture}%
\setlength{\unitlength}{0.00041700in}%
%
\begingroup\makeatletter\ifx\SetFigFont\undefined
% extract first six characters in \fmtname
\def\x#1#2#3#4#5#6#7\relax{\def\x{#1#2#3#4#5#6}}%
\expandafter\x\fmtname xxxxxx\relax \def\y{splain}%
\ifx\x\y   % LaTeX or SliTeX?
\gdef\SetFigFont#1#2#3{%
  \ifnum #1<17\tiny\else \ifnum #1<20\small\else
  \ifnum #1<24\normalsize\else \ifnum #1<29\large\else
  \ifnum #1<34\Large\else \ifnum #1<41\LARGE\else
     \huge\fi\fi\fi\fi\fi\fi
  \csname #3\endcsname}%
\else
\gdef\SetFigFont#1#2#3{\begingroup
  \count@#1\relax \ifnum 25<\count@\count@25\fi
  \def\x{\endgroup\@setsize\SetFigFont{#2pt}}%
  \expandafter\x
    \csname \romannumeral\the\count@ pt\expandafter\endcsname
    \csname @\romannumeral\the\count@ pt\endcsname
  \csname #3\endcsname}%
\fi
\fi\endgroup
\begin{picture}(10342,1339)(130,-1341)
\put(1186,-811){\makebox(0,0)[lb]{\smash{\SetFigFont{12}{14.4}{rm}$=$}}}
\put(4306,-781){\makebox(0,0)[lb]{\smash{\SetFigFont{12}{14.4}{rm}$=$}}}
\put(2791,-781){\makebox(0,0)[lb]{\smash{\SetFigFont{12}{14.4}{rm}$-$}}}
\put(5956,-826){\makebox(0,0)[lb]{\smash{\SetFigFont{12}{14.4}{rm}$-$}}}
\put(9196,-811){\makebox(0,0)[lb]{\smash{\SetFigFont{12}{14.4}{rm}$-$}}}
\put(9991,-751){\makebox(0,0)[lb]{\smash{\SetFigFont{12}{14.4}{rm}$e$}}}
\put(7291,-811){\makebox(0,0)[lb]{\smash{\SetFigFont{12}{14.4}{rm}$+$}}}
\put(7876,-811){\makebox(0,0)[lb]{\smash{\SetFigFont{12}{14.4}{rm}$\frac{1}{2}$}}}
\end{picture}

\]
\ni In general by [LM], $P_L(1,-1)$ is $(-1)^{c-1}(-2)^{\frac
12 d}$ where $c$ is the number of components of $L$ and $d$ is the
dimension of the first homology group (with $\Bbb Z/2\Bbb Z$
coefficients) of the triple branched cover of $S^3$, branched over
$L$. It would be reassuring to be able to see directly why our
formula gives this value. This would also prove directly that the
map $\crsst\mapsto \boxe{e}-\frac{1}{2} \vln\ \ \smth$
passes to the {\sc homfly} quotient.
Our proof of this is a little indirect --- one shows that the
planar subalgebras $P^H$ and $(P^\s)^{A_4}$ are the same by showing
they arise as centralizer towers from the same subfactor
(constructed in [GHJ]). Thus there must be a 3-box
corresponding to: $\crsst$  and we obtained the explicit
expression for it by solving an obvious set of equations.

As far as we know, this is the first genuine ``3-spin interaction"
statistical mechanical model for a link invariant.
Of course one may produce 3-spin interaction models by taking a
2-spin one and summing over the internal spin $\s$ in the picture
\[
        \begin{picture}(0,0)%
\epsfig{file=ps/pic435.ps}%
\end{picture}%
\setlength{\unitlength}{0.00041700in}%
%
\begingroup\makeatletter\ifx\SetFigFont\undefined
% extract first six characters in \fmtname
\def\x#1#2#3#4#5#6#7\relax{\def\x{#1#2#3#4#5#6}}%
\expandafter\x\fmtname xxxxxx\relax \def\y{splain}%
\ifx\x\y   % LaTeX or SliTeX?
\gdef\SetFigFont#1#2#3{%
  \ifnum #1<17\tiny\else \ifnum #1<20\small\else
  \ifnum #1<24\normalsize\else \ifnum #1<29\large\else
  \ifnum #1<34\Large\else \ifnum #1<41\LARGE\else
     \huge\fi\fi\fi\fi\fi\fi
  \csname #3\endcsname}%
\else
\gdef\SetFigFont#1#2#3{\begingroup
  \count@#1\relax \ifnum 25<\count@\count@25\fi
  \def\x{\endgroup\@setsize\SetFigFont{#2pt}}%
  \expandafter\x
    \csname \romannumeral\the\count@ pt\expandafter\endcsname
    \csname @\romannumeral\the\count@ pt\endcsname
  \csname #3\endcsname}%
\fi
\fi\endgroup
\begin{picture}(3344,2444)(1479,-3683)
\put(2326,-1711){\makebox(0,0)[lb]{\smash{\SetFigFont{12}{14.4}{rm}$\sigma_1$}}}
\put(1651,-3436){\makebox(0,0)[lb]{\smash{\SetFigFont{12}{14.4}{rm}$\sigma3$}}}
\put(2626,-2761){\makebox(0,0)[lb]{\smash{\SetFigFont{12}{14.4}{rm}$\sigma$}}}
\put(4426,-3436){\makebox(0,0)[lb]{\smash{\SetFigFont{12}{14.4}{rm}$\sigma_2$}}}
\end{picture}

\]
\ni but that is of little interest. One may check quite easily that
the above model does not factorize in this way.

\bs\ni{\bf Example 2.9. Finite groups}

A special case of the subalgebra of $P^\s$ of Example 2.8 where a
group acts on a set $S$ is where $S=G$ itself, the action being
left multiplication. The algebra $P_2$ is then the group algebra
$\Bbb C G$, linearly spanned by elements $\boxg{g}$ for $g$
in $G$, which are by definition the matrices
$g^a_b=\d_{ag,b}$. They have the properties
%%%%%%%%%%%%%%%%%%%%%%%%%%%%%%%%%%%%%%%%%%%%%%%%
\[
        \begin{picture}(0,0)%
\epsfig{file=ps/pic43.ps}%
\end{picture}%
\setlength{\unitlength}{0.00041700in}%
%
\begingroup\makeatletter\ifx\SetFigFont\undefined
% extract first six characters in \fmtname
\def\x#1#2#3#4#5#6#7\relax{\def\x{#1#2#3#4#5#6}}%
\expandafter\x\fmtname xxxxxx\relax \def\y{splain}%
\ifx\x\y   % LaTeX or SliTeX?
\gdef\SetFigFont#1#2#3{%
  \ifnum #1<17\tiny\else \ifnum #1<20\small\else
  \ifnum #1<24\normalsize\else \ifnum #1<29\large\else
  \ifnum #1<34\Large\else \ifnum #1<41\LARGE\else
     \huge\fi\fi\fi\fi\fi\fi
  \csname #3\endcsname}%
\else
\gdef\SetFigFont#1#2#3{\begingroup
  \count@#1\relax \ifnum 25<\count@\count@25\fi
  \def\x{\endgroup\@setsize\SetFigFont{#2pt}}%
  \expandafter\x
    \csname \romannumeral\the\count@ pt\expandafter\endcsname
    \csname @\romannumeral\the\count@ pt\endcsname
  \csname #3\endcsname}%
\fi
\fi\endgroup
\begin{picture}(8914,1751)(879,-1467)
\put(1051,-1096){\makebox(0,0)[lb]{\smash{\SetFigFont{12}{14.4}{rm}$h$}}}
\put(4336,-676){\makebox(0,0)[lb]{\smash{\SetFigFont{12}{14.4}{rm}$=$}}}
\put(1531,-691){\makebox(0,0)[lb]{\smash{\SetFigFont{12}{14.4}{rm}$=$}}}
\put(2551,-781){\makebox(0,0)[lb]{\smash{\SetFigFont{12}{14.4}{rm},}}}
\put(5986,-781){\makebox(0,0)[lb]{\smash{\SetFigFont{12}{14.4}{rm},}}}
\put(1021,-211){\makebox(0,0)[lb]{\smash{\SetFigFont{12}{14.4}{rm}$g$}}}
\put(1996,-646){\makebox(0,0)[lb]{\smash{\SetFigFont{12}{14.4}{rm}$gh$}}}
\put(3961,-661){\makebox(0,0)[lb]{\smash{\SetFigFont{12}{14.4}{rm}$h$}}}
\put(3106,-646){\makebox(0,0)[lb]{\smash{\SetFigFont{12}{14.4}{rm}$g$}}}
\put(8476,-241){\makebox(0,0)[lb]{\smash{\SetFigFont{12}{14.4}{rm}$g$}}}
\put(6556,-826){\makebox(0,0)[lb]{\smash{\SetFigFont{12}{14.4}{rm}$g$}}}
\put(5446,-646){\makebox(0,0)[lb]{\smash{\SetFigFont{12}{14.4}{rm}$g$}}}
\put(4711,-676){\makebox(0,0)[lb]{\smash{\SetFigFont{12}{14.4}{rm}$\delta_{g,h}$}}}
\put(9376,-961){\makebox(0,0)[lb]{\smash{\special{ps:gsave currentpoint currentpoint translate
-90.0 rotate neg exch neg exch translate}\SetFigFont{12}{14.4}{rm}$gh$\special{ps:currentpoint grestore moveto}}}}
\put(7501,-211){\makebox(0,0)[lb]{\smash{\special{ps:gsave currentpoint currentpoint translate
-90.0 rotate neg exch neg exch translate}\SetFigFont{12}{14.4}{rm}$h$\special{ps:currentpoint grestore moveto}}}}
\end{picture}

\]
\ni which can be used to present $P$ as a planar algebra (see
[La]). In particular $P_{1,3}=\ell^{\i}(G)$ and $P_2=M_{|G|}(\Bbb
C)$. It is obvious that $G$ can be recovered from the abstract
planar algebra $P$ by using the minimal idempotents of $P_{1,3}$.

A noncommutative, finite-dimensional Hopf algebra  gives a planar
algebra but it cannot be a subalgebra of $P^\s$.

\bs{\bf Example 2.10 \ Invariant planar algebras.} Given an
invertible element $u\in P_1$ in a general planar algebra $P$ we
define $u^{\otimes k}$ to be the element of $P_k$ defined by the
following $k$-tangle
\[
        \begin{picture}(0,0)%
\epsfig{file=ps/pic64.ps}%
\end{picture}%
\setlength{\unitlength}{0.00041700in}%
%
\begingroup\makeatletter\ifx\SetFigFont\undefined
% extract first six characters in \fmtname
\def\x#1#2#3#4#5#6#7\relax{\def\x{#1#2#3#4#5#6}}%
\expandafter\x\fmtname xxxxxx\relax \def\y{splain}%
\ifx\x\y   % LaTeX or SliTeX?
\gdef\SetFigFont#1#2#3{%
  \ifnum #1<17\tiny\else \ifnum #1<20\small\else
  \ifnum #1<24\normalsize\else \ifnum #1<29\large\else
  \ifnum #1<34\Large\else \ifnum #1<41\LARGE\else
     \huge\fi\fi\fi\fi\fi\fi
  \csname #3\endcsname}%
\else
\gdef\SetFigFont#1#2#3{\begingroup
  \count@#1\relax \ifnum 25<\count@\count@25\fi
  \def\x{\endgroup\@setsize\SetFigFont{#2pt}}%
  \expandafter\x
    \csname \romannumeral\the\count@ pt\expandafter\endcsname
    \csname @\romannumeral\the\count@ pt\endcsname
  \csname #3\endcsname}%
\fi
\fi\endgroup
\begin{picture}(5326,1094)(1951,-2932)
\put(3708,-2446){\makebox(0,0)[lb]{\smash{\SetFigFont{12}{14.4}{rm}$u$}}}
\put(5508,-2459){\makebox(0,0)[lb]{\smash{\SetFigFont{12}{14.4}{rm}$u$}}}
\put(4951,-2228){\makebox(0,0)[lb]{\smash{\special{ps:gsave currentpoint currentpoint translate
-180.0 rotate neg exch neg exch translate}\SetFigFont{12}{14.4}{rm}$u$\special{ps:currentpoint grestore moveto}}}}
\put(6708,-2459){\makebox(0,0)[lb]{\smash{\SetFigFont{12}{14.4}{rm}$u$}}}
\put(5957,-2505){\makebox(0,0)[lb]{\smash{\SetFigFont{12}{14.4}{rm}$\cdots$}}}
\put(1951,-2536){\makebox(0,0)[lb]{\smash{\SetFigFont{12}{14.4}{rm}$u^{\otimes k}=$}}}
\end{picture}

\]
%%%%%%%%%%%%%%%%%%%%%%%%%%%%%%%%%%%%%%%%%%%%%%%
%$$
%u^{\otimes k} = \boxed{ \ \boxed{u} \ \boxed{n} \
% \boxed{u} \ \dots  \boxed{u} \  }
%$$
%%%%%%%%%%%%%%%%%%%%%%%%%%%%%%%%
($k$ is odd in the picture). The $u$'s and upside down $u^{-1}$'s alternate.

\medskip
{\bf Proposition 2.10.1} If $P$ is a general planar algebra
and $S$ is a set of invertible elements of $P_1$, set
$$
P^S_k=\{x\in P_k\mid u^{\otimes k}x=xu^{\otimes k} \ \forall \
u\in S\} \ .
$$
Then $P^S$ is a general planar subalgebra of $P$ (planar if $P$ is)
and a $*$-planar algebra if $P$ is, and $S=S^*$.

\bigskip
{\sc Proof.} That $P^S$ is a unital filtered subalgebra is
obvious. In the $*$ case note that
$(u^{\otimes k})^*=(u^*)^{\otimes k}$. So by Lemma 1.18 we only
have to check invariance under ${\cal A}(\phi)$. Given an
$A\in{\cal A}_{k,n}(\phi)$ consider the tangle representing
$(\otimes^nu)\pi_A(x)(\otimes ^nu)^{-1}$ in Figure~\ref{pic66}

\begin{figure}[h]
\[
        \begin{picture}(0,0)%
\epsfig{file=ps/pic66.ps}%
\end{picture}%
\setlength{\unitlength}{0.00041700in}%
%
\begingroup\makeatletter\ifx\SetFigFont\undefined
% extract first six characters in \fmtname
\def\x#1#2#3#4#5#6#7\relax{\def\x{#1#2#3#4#5#6}}%
\expandafter\x\fmtname xxxxxx\relax \def\y{splain}%
\ifx\x\y   % LaTeX or SliTeX?
\gdef\SetFigFont#1#2#3{%
  \ifnum #1<17\tiny\else \ifnum #1<20\small\else
  \ifnum #1<24\normalsize\else \ifnum #1<29\large\else
  \ifnum #1<34\Large\else \ifnum #1<41\LARGE\else
     \huge\fi\fi\fi\fi\fi\fi
  \csname #3\endcsname}%
\else
\gdef\SetFigFont#1#2#3{\begingroup
  \count@#1\relax \ifnum 25<\count@\count@25\fi
  \def\x{\endgroup\@setsize\SetFigFont{#2pt}}%
  \expandafter\x
    \csname \romannumeral\the\count@ pt\expandafter\endcsname
    \csname @\romannumeral\the\count@ pt\endcsname
  \csname #3\endcsname}%
\fi
\fi\endgroup
\begin{picture}(7994,5444)(1179,-5783)
\put(2161,-1366){\makebox(0,0)[lb]{\smash{\SetFigFont{12}{14.4}{rm}$u$}}}
\put(4576,-1336){\makebox(0,0)[lb]{\smash{\SetFigFont{12}{14.4}{rm}$u$}}}
\put(6991,-1336){\makebox(0,0)[lb]{\smash{\SetFigFont{12}{14.4}{rm}$u$}}}
\put(4366,-4876){\makebox(0,0)[lb]{\smash{\SetFigFont{12}{14.4}{rm}$u^{-1}$}}}
\put(6736,-4876){\makebox(0,0)[lb]{\smash{\SetFigFont{12}{14.4}{rm}$u^{-1}$}}}
\put(3721,-1171){\makebox(0,0)[lb]{\smash{\special{ps:gsave currentpoint currentpoint translate
-180.0 rotate neg exch neg exch translate}\SetFigFont{12}{14.4}{rm}$u^{-1}$\special{ps:currentpoint grestore moveto}}}}
\put(6091,-1171){\makebox(0,0)[lb]{\smash{\special{ps:gsave currentpoint currentpoint translate
-180.0 rotate neg exch neg exch translate}\SetFigFont{12}{14.4}{rm}$u^{-1}$\special{ps:currentpoint grestore moveto}}}}
\put(8536,-1186){\makebox(0,0)[lb]{\smash{\special{ps:gsave currentpoint currentpoint translate
-180.0 rotate neg exch neg exch translate}\SetFigFont{12}{14.4}{rm}$u^{-1}$\special{ps:currentpoint grestore moveto}}}}
\put(3691,-4621){\makebox(0,0)[lb]{\smash{\special{ps:gsave currentpoint currentpoint translate
-180.0 rotate neg exch neg exch translate}\SetFigFont{12}{14.4}{rm}$u$\special{ps:currentpoint grestore moveto}}}}
\put(6076,-4636){\makebox(0,0)[lb]{\smash{\special{ps:gsave currentpoint currentpoint translate
-180.0 rotate neg exch neg exch translate}\SetFigFont{12}{14.4}{rm}$u$\special{ps:currentpoint grestore moveto}}}}
\put(8416,-4651){\makebox(0,0)[lb]{\smash{\special{ps:gsave currentpoint currentpoint translate
-180.0 rotate neg exch neg exch translate}\SetFigFont{12}{14.4}{rm}$u$\special{ps:currentpoint grestore moveto}}}}
\put(1936,-4861){\makebox(0,0)[lb]{\smash{\SetFigFont{12}{14.4}{rm}$u^{-1}$}}}
\end{picture}

\]
\caption[2.10.2]{}\label{pic66}
\end{figure}
\ni Each string of $\otimes^ku$, at the top and bottom, either
connects to another external boundary point, in which case the $u$
cancels with $u^{-1}$, or it connects to an internal boundary
point of the annulus. Isotoping each such $u$ and $u^{-1}$ close to
the internal boundary and inserting cancelling pairs of $u$ and
$u^{-1}$ on strings connecting internal boundary points, we see
that the tangle of Figure 2.10.2 gives the same element of $P$ as
the one where the only instances of $u$ and $u^{-1}$ surround the
$\blankbox{x}$ in an alternating fashion.
Since $x\in P^S$, these $u$'s may be eliminated and we are left
with $\pi_A(x)$. \qed

\bs
This gives a useful way of constructing planar subalgebras. In
particular $P^S_1$ is the commutant of $S$ in $P_1$ which may be
much smaller than $P_1$. Of special interest is the case where
$P=P^{\otimes}$ and $S$ is a subgroup $G$ of the unitary group.
In this case $P^G_k={\op{End}}_G(V\otimes V^*\otimes V\otimes
V^*\otimes \dots)$ ($k$ copies of $V$ or $V^*$) where $V^*$
represents the contragredient representation of $G$. Other cases of
interest can be constructed by cabling as in $\S$3 and then picking
some set of invertible elements in the original  $P_n$.

To obtain a more general construction one may replace
$\otimes^nu$ with tangles of the form
\[
        \begin{picture}(0,0)%
\epsfig{file=ps/pic67.ps}%
\end{picture}%
\setlength{\unitlength}{0.00041700in}%
%
\begingroup\makeatletter\ifx\SetFigFont\undefined
% extract first six characters in \fmtname
\def\x#1#2#3#4#5#6#7\relax{\def\x{#1#2#3#4#5#6}}%
\expandafter\x\fmtname xxxxxx\relax \def\y{splain}%
\ifx\x\y   % LaTeX or SliTeX?
\gdef\SetFigFont#1#2#3{%
  \ifnum #1<17\tiny\else \ifnum #1<20\small\else
  \ifnum #1<24\normalsize\else \ifnum #1<29\large\else
  \ifnum #1<34\Large\else \ifnum #1<41\LARGE\else
     \huge\fi\fi\fi\fi\fi\fi
  \csname #3\endcsname}%
\else
\gdef\SetFigFont#1#2#3{\begingroup
  \count@#1\relax \ifnum 25<\count@\count@25\fi
  \def\x{\endgroup\@setsize\SetFigFont{#2pt}}%
  \expandafter\x
    \csname \romannumeral\the\count@ pt\expandafter\endcsname
    \csname @\romannumeral\the\count@ pt\endcsname
  \csname #3\endcsname}%
\fi
\fi\endgroup
\begin{picture}(5638,1222)(2004,-1261)
\put(3076,-736){\makebox(0,0)[lb]{\smash{\SetFigFont{12}{14.4}{rm}$u_1$}}}
\put(4651,-736){\makebox(0,0)[lb]{\smash{\SetFigFont{12}{14.4}{rm}$u_2$}}}
\put(6226,-736){\makebox(0,0)[lb]{\smash{\SetFigFont{12}{14.4}{rm}$u_3$}}}
\end{picture}

\]
\ni for $u$'s satisfying appropriate equations. We will use this
approach with $n=1$ to pick up some important cases of commuting
squares in Example 2.11.



\newpage
\ni{\bf 2.11. Binunitaries}

Commuting squares have been used to construct subfactors (see
[GHJ],[Ha]) and the general theory of calculating the
subfactor planar algebra of $\S$4.2 from a commuting square will be
dealt with in a future paper. The treatment uses the language of
statistical mechanical models with some attention paid to critical
points as in chapter 4 of [JS]. Here we give a different
approach which seems more natural from a planar point of view and
will allow us to capture, as special cases, spin model commuting
squares and some vertex model ones with no extension of the planar
algebra formalism. The main concept is that of a bi-invertible
element in a planar algebra which is the next step in the
hierarchy discussed at the end of Example 2.10.

\bs{\bf Definition 2.11.1.} Let $P$ be a general planar algebra. An
invertible element $u\in P_2$ will be called {\it bi-invertible} if
\[
        \begin{picture}(0,0)%
\epsfig{file=ps/pic44.ps}%
\end{picture}%
\setlength{\unitlength}{0.00041700in}%
%
\begingroup\makeatletter\ifx\SetFigFont\undefined
% extract first six characters in \fmtname
\def\x#1#2#3#4#5#6#7\relax{\def\x{#1#2#3#4#5#6}}%
\expandafter\x\fmtname xxxxxx\relax \def\y{splain}%
\ifx\x\y   % LaTeX or SliTeX?
\gdef\SetFigFont#1#2#3{%
  \ifnum #1<17\tiny\else \ifnum #1<20\small\else
  \ifnum #1<24\normalsize\else \ifnum #1<29\large\else
  \ifnum #1<34\Large\else \ifnum #1<41\LARGE\else
     \huge\fi\fi\fi\fi\fi\fi
  \csname #3\endcsname}%
\else
\gdef\SetFigFont#1#2#3{\begingroup
  \count@#1\relax \ifnum 25<\count@\count@25\fi
  \def\x{\endgroup\@setsize\SetFigFont{#2pt}}%
  \expandafter\x
    \csname \romannumeral\the\count@ pt\expandafter\endcsname
    \csname @\romannumeral\the\count@ pt\endcsname
  \csname #3\endcsname}%
\fi
\fi\endgroup
\begin{picture}(8744,1558)(429,-1584)
\put(7276,-586){\makebox(0,0)[lb]{\smash{\special{ps:gsave currentpoint currentpoint translate
-180.0 rotate neg exch neg exch translate}\SetFigFont{12}{14.4}{rm}$u^{-1}$\special{ps:currentpoint grestore moveto}}}}
\put(901,-811){\makebox(0,0)[lb]{\smash{\SetFigFont{12}{14.4}{rm}$u$}}}
\put(3076,-586){\makebox(0,0)[lb]{\smash{\special{ps:gsave currentpoint currentpoint translate
-180.0 rotate neg exch neg exch translate}\SetFigFont{12}{14.4}{rm}$u^{-1}$\special{ps:currentpoint grestore moveto}}}}
\put(3451,-886){\makebox(0,0)[lb]{\smash{\SetFigFont{12}{14.4}{rm}$=$}}}
\put(5776,-886){\makebox(0,0)[lb]{\smash{\SetFigFont{12}{14.4}{rm}$=$}}}
\put(8626,-811){\makebox(0,0)[lb]{\smash{\SetFigFont{12}{14.4}{rm}$u$}}}
\put(4051,-886){\makebox(0,0)[lb]{\smash{\SetFigFont{12}{14.4}{rm}$\Delta$}}}
\end{picture}

\]
\ni for some non-zero scalar $\D$. (If $P$ is planar $\D$ is
necessarily $\d_1 \ (\d_2)$. If $P$ is a planar $*$-algebra, a
 bi-invertible $u$ is called biunitary  if $u^*=u^{-1}$.
Bi-invertible elements define planar subalgebras as we now
describe. It will be convenient to consider labelled tangles
containing certain distinguished curves joining boundary points,
which intersect with only the other strings in a tangle and do not
meet any internal boxes. From such a tangle, and a
bi-invertible element $u$, we construct an honest labelled tangle,
in the sense of $\S$1, in two steps.


